
\section{IMÁGENES} \label{sec:imagenes}

Las imágenes se insertan mediante el comando \textbackslash\texttt{includegraphics} que es conveniente situar en el entorno \texttt{figure} (mismo entorno utilizado, como se verá más adelante, para gráficas o diagramas). Un ejemplo de imagen se muestra a continuación (al utilizar Overleaf es esencial cargar la imagen en el directorio de trabajo antes de insertarla en el documento): 

\vspace{10pt}


En la inmensa mayoría de casos el tamaño original de la imagen no se adapta correctamente a las dimensiones de la página, por lo que es necesario redimensionar la imagen mediante el argumento \texttt{scale} de \textbackslash\texttt{includegraphics}. 

Al igual que para las tablas, existen distintas alternativas en cuanto a su posicionamiento (que, se recuerda, pueden consultarse en \url{https://www.overleaf.com/learn/latex/positioning_images_and_tables}), el título se indica mediante el comando \textbackslash\texttt{caption} y la referencia mediante el comando \textbackslash\texttt{label}. 

Existen, además, diversas opciones en lo relativo al manejo de imágenes que no se detallan en este documento, pero que pueden consultarse en \url{https://es.overleaf.com/learn/latex/Inserting_Images} 
